\chapter{Diffusion Models: Mathematical Foundations}

\section{Introduction}

We now formalize diffusion models as probabilistic processes. The key idea is to define a forward Markov chain that gradually corrupts data into noise, and then learn a reverse process that reconstructs data.

This chapter derives the forward and reverse processes, shows how training objectives arise, and connects diffusion to score matching and stochastic differential equations.

\section{Discrete-Time Diffusion Formulation}

We define a forward Markov chain:
\[
q(x_{1:T} \mid x_0) = \prod_{t=1}^T q(x_t \mid x_{t-1}),
\]
where each transition adds Gaussian noise.

\subsection{Forward Transitions}

\begin{definition}[Forward Process]
\[
q(x_t \mid x_{t-1})
=
\mathcal{N}\left(x_t \;\middle|\; \sqrt{1-\beta_t}\,x_{t-1}, \;\beta_t I \right),
\]
where $\beta_t \in (0,1)$ is a noise schedule.
\end{definition}

This defines a Markov chain that gradually destroys structure.

\section{Forward Process: Closed-Form Marginals}

A key property is that $x_t$ can be sampled directly from $x_0$.

\begin{proposition}[Closed-Form Marginal]
\[
q(x_t \mid x_0)
=
\mathcal{N}\left(x_t \;\middle|\; \sqrt{\bar{\alpha}_t}\,x_0, \;(1-\bar{\alpha}_t)I \right),
\]
where:
\[
\alpha_t = 1 - \beta_t,
\quad
\bar{\alpha}_t = \prod_{s=1}^t \alpha_s.
\]
\end{proposition}

\subsection{Derivation Sketch}

Repeated substitution gives:
\[
x_t = \sqrt{\bar{\alpha}_t} x_0 + \sqrt{1-\bar{\alpha}_t}\,\epsilon,
\quad \epsilon \sim \mathcal{N}(0,I).
\]

Thus:
\begin{itemize}
    \item mean shrinks toward zero,
    \item variance increases over time.
\end{itemize}

\section{Reverse Process Parameterization}

We define a learned reverse chain:
\[
p_\theta(x_{0:T}) = p(x_T)\prod_{t=1}^T p_\theta(x_{t-1} \mid x_t),
\]
with:
\[
p(x_T) = \mathcal{N}(0,I).
\]

\subsection{Gaussian Parameterization}

\[
p_\theta(x_{t-1} \mid x_t)
=
\mathcal{N}(x_{t-1} \mid \mu_\theta(x_t,t), \Sigma_\theta(x_t,t)).
\]

The model learns to predict the mean (and optionally variance) of the reverse transition.

\section{Training Objectives: ELBO View}

Training maximizes log-likelihood:
\[
\log p_\theta(x_0).
\]

Using variational inference, we obtain an ELBO:
\[
\mathbb{E}_q \left[
\log p_\theta(x_{0:T}) - \log q(x_{1:T} \mid x_0)
\right].
\]

This decomposes into terms comparing forward and reverse transitions.

\subsection{Simplified Objective}

After algebraic simplification, training reduces to:
\[
\mathbb{E}_{t,x_0,\epsilon}
\left[
\| \epsilon - \epsilon_\theta(x_t,t) \|^2
\right],
\]
where:
\[
x_t = \sqrt{\bar{\alpha}_t}x_0 + \sqrt{1-\bar{\alpha}_t}\epsilon.
\]

\subsection{Interpretation}

\begin{itemize}
    \item model predicts noise $\epsilon$,
    \item equivalent to learning reverse dynamics,
    \item yields stable supervised learning objective.
\end{itemize}

\section{Training Objective Transformations}

Different parameterizations are equivalent:

\begin{itemize}
    \item predicting noise $\epsilon$,
    \item predicting $x_0$,
    \item predicting score $\nabla_x \log p(x_t)$.
\end{itemize}

These correspond to different but related loss formulations.

\section{Score Matching and SDE View}

Diffusion models are closely related to score matching.

\subsection{Score Function}

\[
s(x) = \nabla_x \log p(x).
\]

The model learns an approximation:
\[
s_\theta(x_t, t).
\]

\subsection{Connection}

Noise prediction is proportional to the score:
\[
\epsilon_\theta(x_t,t) \;\Longleftrightarrow\; -\nabla_x \log q(x_t).
\]

Thus diffusion training implicitly performs score matching.

\subsection{Continuous-Time View}

As $T \to \infty$, the process becomes a stochastic differential equation (SDE):
\[
dx = f(x,t)\,dt + g(t)\,dW_t.
\]

The reverse process is another SDE that removes noise.

\section{Langevin Dynamics Connection}

Sampling can be interpreted as approximate Langevin dynamics:
\[
x_{t-1} = x_t + \eta \nabla_x \log p(x_t) + \sqrt{2\eta}\,\xi.
\]

\begin{itemize}
    \item gradient term moves toward high-density regions,
    \item noise maintains stochasticity,
    \item diffusion combines both ideas.
\end{itemize}

\section{DDPM, DDIM, and Sampler Intuition}

\subsection{DDPM (Stochastic Sampling)}

\begin{itemize}
    \item follows learned reverse transitions,
    \item includes stochastic noise at each step,
    \item produces diverse samples.
\end{itemize}

\subsection{DDIM (Deterministic Sampling)}

\begin{itemize}
    \item removes stochasticity,
    \item produces deterministic trajectories,
    \item faster sampling with fewer steps.
\end{itemize}

\subsection{Tradeoffs}

\begin{itemize}
    \item stochastic samplers: more diversity,
    \item deterministic samplers: faster and more controllable.
\end{itemize}

\section{A Unifying Perspective}

Diffusion models define:
\begin{itemize}
    \item a forward stochastic process that simplifies data,
    \item a reverse learned process that reconstructs it.
\end{itemize}

Training can be viewed as:
\begin{itemize}
    \item variational inference (ELBO),
    \item score matching,
    \item learning dynamics of a stochastic system.
\end{itemize}

Thus diffusion sits at the intersection of probability, optimization, and dynamical systems.

\section{Summary}

In this chapter we formalized diffusion models. We defined the forward Markov process, derived closed-form marginals, and introduced the reverse generative model. We showed how training arises from a variational objective and reduces to noise prediction.

We also connected diffusion to score matching, stochastic differential equations, and Langevin dynamics, and discussed different sampling schemes such as DDPM and DDIM.

\section*{Exercises}

\begin{enumerate}
    \item Derive the mean and variance of $q(x_t \mid x_0)$.

    \item Show how $x_t$ can be written as a linear combination of $x_0$ and Gaussian noise.

    \item Derive the simplified noise-prediction objective from the ELBO.

    \item Explain why predicting $\epsilon$ is equivalent to learning the score function.

    \item Compare stochastic (DDPM) and deterministic (DDIM) sampling.

    \item Show how variance accumulates over diffusion steps.

    \item Derive the relationship between score matching and denoising.

    \item Explain the role of the noise schedule $\{\beta_t\}$.
\end{enumerate}